\medskip

\textbf{Partie A}

\medskip

On considère l'équation différentielle

\[(E)\: \qquad y' + 0,4y = \e^{-0,4t}\]

où $y$ est une fonction de la variable réelle $t$.

On cherche l'ensemble des fonctions définies et dérivables sur $\R$ qui sont solutions
de cette équation.

\medskip

\begin{enumerate}
\item Soit $u$ la fonction définie sur $\R$ par : $u(t) = t\e^{-0,4t}$.

Vérifier que $u$ est solution de $(E)$.
\item Soit $f$ une fonction définie et dérivable sur $\R$.

On note $g$ la fonction définie sur $\R$ par : $g(t) = f(t) - u(t)$.

Soit $(H)$ l'équation différentielle $y' + 0,4y = 0$.
	\begin{enumerate}
		\item Démontrer que si la fonction $g$ est solution de l'équation différentielle $(H)$ alors la fonction $f$ est solution de l'équation différentielle $(E)$.
	\end{enumerate}		
		
On admettra que la réciproque est vraie.

	\begin{enumerate}[resume]
		\item Résoudre l'équation différentielle $(H)$.
		\item En déduire les solutions de $(E)$.
		\item Déterminer la solution $f$ de $(E)$ telle que $f(0) = 1$.
	\end{enumerate}
\end{enumerate}

\bigskip

\textbf{Partie B}

\medskip

On s'intéresse à la glycémie chez une personne venant de prendre un repas.

La glycémie en g.L$^{-1}$, en fonction du temps $t$, exprimé en heure, écoulé depuis la fin
du repas, est modélisée par la fonction $f$ définie sur [0~;~6] par :

\[f(t) = (t + 1)\e^{-0,4t}.\]

\medskip

\begin{enumerate}
\item
	\begin{enumerate}
		\item Montrer que, pour tout $t \in [0~;~6],\: f'(t) = (- 0,4t + 0,6)\e^{-0,4t}$.
		\item Étudier les variations de $f$ sur [0~;~6] puis dresser son tableau de variations sur cet intervalle.
	\end{enumerate}
\item Une personne est en hypoglycémie lorsque sa glycémie est inférieure à 0,7~g.L$^{-1}$.
	\begin{enumerate}
		\item Démontrer que sur l'intervalle [0~;~6] l'équation $f(t) = 0,7$ admet une unique
solution que l'on notera $\alpha$.
		\item Au bout de combien de temps après avoir pris son repas cette personne est-
elle en hypoglycémie ?

On exprimera ce temps à la minute près.
	\end{enumerate}
\item On souhaite déterminer la glycémie moyenne en g.L$^{-1}$ chez cette personne lors des six heures qui suivent le repas.
	\begin{enumerate}
		\item À l'aide d'une intégration par parties, montrer que:

\[\displaystyle\int_0^6 f(t)\: \text{d}t = - 23,75\e^{-2,4} + 8,75.\]

		\item Calculer la glycémie moyenne en g.L$^{-1}$ chez cette personne lors des six heures qui suivent le repas.
		\item En remarquant que la fonction $f$ est solution de l'équation différentielle $(E)$, expliquer comment on aurait pu obtenir ce résultat autrement.
	\end{enumerate}
\end{enumerate}

\bigskip
